\section{Single-flywheel launcher: exit speed and spin}
\label{sec:launcher}

% Source of truth: src/launcher.py (simulate_shot, hood_ideal, brettle_ratio, brettle_omega_f,
% frictionless_top_limit, roller_rolling_state, roller_coast, recovery_time) and
% src/shooter_spec.py (wheel_rpm, spin_rpm, V_OVER_U). Numbers below were recomputed by running
% that code with the system python on 2026-10-30; data/launch_envelope.json holds the stored envelope.

The flight sections take the exit state of the ball as an input: an exit speed $v$ and a spin
parameter $S$. This section asks the opposite question. Given \emph{one} powered 98\,\unit{mm}\userval{}
flywheel with either a fixed curved hood or one passive (unpowered) upper roller, which pairs
$(v,S)$ can the launcher actually produce, how fast must the wheel turn, and what does one shot cost
the wheel? Everything here follows \code{src/launcher.py}. The model is uncalibrated: no POLLEN has been
measured on a launcher yet, so every simulated output is a prediction.

\subsection{Frames, sign conventions and the wheel surface speed}
\label{sec:launcher:frames}

The model is planar and one-dimensional along the contact path. Let $x$ be the arc length travelled by
the ball centre along the path (\unit{m}), positive in the direction of travel, and let $z$ point away from
the wheel. The state is
\begin{itemize}
  \item $v$: speed of the ball centre along the path (\unit{m/s}), positive forward;
  \item $s$: absolute angular velocity of the ball (\unit{rad/s}), positive for \emph{backspin}, i.e.\ the
        wheel-side point of the ball moves forward relative to the centre and the hood-side point moves
        backward;
  \item $\Omega$: flywheel angular velocity (\unit{rad/s}), positive when the wheel surface at the contact
        moves forward;
  \item $\Omega_t$: passive-roller angular velocity (\unit{rad/s}), positive when the roller surface at the
        contact moves forward (a ``counter-rotating'' roller).
\end{itemize}
Wheel speeds enter and leave the code in revolutions per minute, $n$ (\unit{rev/min}), and are
converted with
\begin{equation}
  \Omega = \frac{2\pi}{60}\,n .
  \label{eq:launcher-rpm}
\end{equation}
The wheel has radius $R_w$ (\unit{m}) and diameter $D = 2R_w$. A point on its rim moves at
\begin{equation}
  u = R_w\,\Omega = \frac{D}{2}\cdot\frac{2\pi n}{60} = \frac{\pi D\, n}{60},
  \label{eq:launcher-u}
\end{equation}
which is the \emph{wheel surface speed} (\unit{m/s}). The spin parameter used by every other section is
defined with the \emph{nominal} ball radius $R_b$ (\unit{m}),
\begin{equation}
  S = \frac{s\,R_b}{v},
  \label{eq:launcher-S}
\end{equation}
so $S=1$ means the ball's nominal surface turns as fast as its centre moves.

\paragraph{Worked example.} With $D = 0.098$\,\unit{m}\userval{} and $n = 3000$\,\unit{rev/min}\assumed,
\eqref{eq:launcher-rpm} gives $\Omega = 314.16$\,\unit{rad/s}\derived{} and \eqref{eq:launcher-u} gives
$u = \pi(0.098)(3000)/60 = 15.39$\,\unit{m/s}\derived.

\paragraph{Parameters.} Table~\ref{tab:launcher-params} lists the defaults of \code{PARAMS} and of the
dataclasses \code{Ball}, \code{Wheel}, \code{Contact}, \code{Hood} and \code{Roller}.

\begin{table}[ht]
\centering
\small
\begin{tabular}{@{}llll@{}}
\toprule
Symbol & Meaning & Default & Range used \\
\midrule
$m$ & ball mass & 0.025\,\unit{kg}\confirmed & 0.0245--0.0255\,\unit{kg}\assumed \\
$R_b$ & nominal ball radius & 0.0355\,\unit{m}\confirmed & -- \\
$k_I$ & ball inertia factor, $I_b = k_I m R_b^2$ & 0.63\assumed & 0.55--0.667\assumed \\
$R_w$ & flywheel radius & 0.049\,\unit{m}\userval & -- \\
$I_w$ & flywheel + steel flywheel inertia & $2.65\times10^{-4}$\,\unit{kg\,m$^2$}\assumed & $0.5$--$5\times10^{-4}$\,\unit{kg\,m$^2$}\assumed \\
$\mu_w$ & friction, wheel tread on ball & 0.8\assumed & 0.5--1.2\assumed \\
$\mu_t$ & friction, hood or roller on ball & 0.3\assumed{} (hood), 0.8\assumed{} (roller) & 0.05--1.0\assumed \\
$\delta$ & diametral compression at the gap & 0.004\,\unit{m}\userval & 0.002--0.008\,\unit{m}\assumed \\
$k_n$ & linear diametral stiffness, $N = k_n\delta$ & 4000\,\unit{N/m}\assumed & 1500--10000\,\unit{N/m}\assumed \\
$\lambda$ & share of $\delta$ that shrinks the rolling radius & 1\assumed & 0--1\assumed \\
$c_{rr}$ & rolling-hysteresis lever arm / $r_c$ & 0.02\assumed & 0--0.06\assumed \\
$\theta$ & hood wrap angle & $90^\circ$\assumed & $45$--$150^\circ$\assumed \\
$v_{in}$ & feed speed into the contact (no spin) & 0.5\,\unit{m/s}\assumed & 0.2--1.5\,\unit{m/s}\assumed \\
$R_t$ & roller radius & 0.024\,\unit{m}\assumed & 0.0125--0.050\,\unit{m}\assumed \\
$I_t$ & roller inertia & $1\times10^{-5}$\,\unit{kg\,m$^2$}\assumed & $10^{-6}$--$2\times10^{-4}$\,\unit{kg\,m$^2$}\assumed \\
$\tau_c$ & roller bearing Coulomb torque & $5\times10^{-4}$\,\unit{N\,m}\assumed & 0--$3\times10^{-3}$\,\unit{N\,m}\assumed \\
$c_v$ & roller viscous drag coefficient & $1\times10^{-6}$\,\unit{N\,m\,s}\assumed & 0--$10^{-5}$\,\unit{N\,m\,s}\assumed \\
$\Delta t_{shot}$ & time between shots & 0.5\,\unit{s}\assumed & 0.25--1.5\,\unit{s}\assumed \\
$n_{free}$, $\tau_s$ & goBILDA 5203 1:1 free speed, stall torque & 6000\,\unit{rev/min}\confirmed, 0.144\,\unit{N\,m}\confirmed & -- \\
\bottomrule
\end{tabular}
\caption{Launcher inputs and their provenance (\code{src/launcher.py}, \code{PARAMS}).}
\label{tab:launcher-params}
\end{table}

\subsection{Contact kinematics: where the ball touches and how fast that point moves}
\label{sec:launcher:kinematics}

\paragraph{Motivation.} Friction acts on the \emph{surface} of the ball, so everything depends on how fast
the two touching points move compared with the wheel and the hood.

\paragraph{Ball inertia.} The ball is a hollow perforated shell. The code writes its moment of inertia
about a diameter as
\begin{equation}
  I_b = k_I\, m\, R_b^2 ,
  \label{eq:launcher-Ib}
\end{equation}
with $k_I$ dimensionless ($k_I = 2/3$ for a thin shell, $2/5$ for a solid sphere) and $I_b$ in
\unit{kg\,m$^2$}.

\paragraph{Effective rolling radius.} The ball is squeezed by $\delta$ (diametral, \unit{m}) between the
wheel and the hood. The code assumes a fraction $\lambda$ of that squeeze shortens the lever arm from the
ball centre to each contact, shared equally between the two contacts:
\begin{equation}
  r_c = R_b - \lambda\,\frac{\delta}{2}.
  \label{eq:launcher-rc}
\end{equation}
$\lambda = 1$ puts all of the squeeze in the ball; $\lambda = 0$ is rigid-radius rolling on a flattened
patch.

\paragraph{Velocity of the two contact points.} Use a right-handed frame $(\hat{x},\hat{y},\hat{z})$ with
$\hat{x}$ forward and $\hat{z}$ away from the wheel. With the sign convention above, backspin $s>0$ is the
angular-velocity vector $\vect{\omega} = -s\,\hat{y}$. The wheel contact sits at
$\vect{r}_w = -r_c\hat{z}$ from the ball centre and the hood (or roller) contact at
$\vect{r}_t = +r_c\hat{z}$. Rigid-body kinematics, $\vect{v}_P = \vect{v}_c + \vect{\omega}\times\vect{r}_P$,
gives
\begin{align}
  \vect{v}_{P,w} &= v\hat{x} + (-s\hat{y})\times(-r_c\hat{z}) = (v + s\,r_c)\,\hat{x}, \label{eq:launcher-vpw}\\
  \vect{v}_{P,t} &= v\hat{x} + (-s\hat{y})\times(+r_c\hat{z}) = (v - s\,r_c)\,\hat{x}, \label{eq:launcher-vpt}
\end{align}
using $\hat{y}\times\hat{z} = \hat{x}$. For a hood concentric with the wheel both contact normals pass
through the ball centre and the wheel axis, so the same two expressions hold along the curved path, with
$x$ measured along the arc.

\paragraph{Slip velocities.} The code defines slip as \emph{surface speed minus ball-point speed}:
\begin{align}
  \sigma_w &= R_w\Omega - (v + s\,r_c), \label{eq:launcher-slipw}\\
  \sigma_t &= R_t\Omega_t - (v - s\,r_c), \qquad (\text{hood: } R_t\Omega_t \equiv 0). \label{eq:launcher-slipt}
\end{align}
$\sigma > 0$ means the surface runs ahead of the ball, so its friction pushes the ball \emph{forward}.

\paragraph{Worked example.} With $R_b = 0.0355$\,\unit{m}\confirmed, $\delta = 0.004$\,\unit{m}\userval{} and
$\lambda = 1$\assumed, \eqref{eq:launcher-rc} gives $r_c = 0.0355 - 0.002 = 0.0335$\,\unit{m}\derived. With
$m = 0.025$\,\unit{kg}\confirmed{} and $k_I = 0.63$\assumed, \eqref{eq:launcher-Ib} gives
$I_b = 0.63\times0.025\times0.0355^2 = 1.985\times10^{-5}$\,\unit{kg\,m$^2$}\derived.

\subsection{The ideal hood and where ``$v = u/2$'' comes from}
\label{sec:launcher:ideal}

\paragraph{Motivation.} The most quoted rule of thumb for a hooded flywheel is that the ball leaves at
half the wheel surface speed. This subsection derives it, so that the conditions it needs are visible.

Assume rolling without slip at \emph{both} contacts at the moment of exit, a fixed hood, and a wheel speed
$\Omega_f$ at exit. Setting $\sigma_t = 0$ and $\sigma_w = 0$ in
\eqref{eq:launcher-slipw}--\eqref{eq:launcher-slipt}:
\begin{align}
  0 &= v - s\,r_c, \label{eq:launcher-rollhood}\\
  R_w\Omega_f &= v + s\,r_c. \label{eq:launcher-rollwheel}
\end{align}
Adding the two equations eliminates $s$: $R_w\Omega_f = 2v$. Subtracting them eliminates $v$:
$R_w\Omega_f = 2 s r_c$. Hence
\begin{equation}
  v = \frac{R_w\,\Omega_f}{2}, \qquad s = \frac{v}{r_c}, \qquad
  S = \frac{s R_b}{v} = \frac{R_b}{r_c}.
  \label{eq:launcher-ideal}
\end{equation}
Three things are hidden in \eqref{eq:launcher-ideal}. (i) $\Omega_f$ is the wheel speed \emph{at exit}, which
is lower than the speed before the shot (Section~\ref{sec:launcher:drop}). (ii) Both contacts must have
stopped slipping, which needs enough friction and enough contact length
(Section~\ref{sec:launcher:grip}). (iii) The compression does not change $v$ but does raise $S$ above 1.
The code function \code{hood\_ideal} uses the rigid case $r_c = R_b$ (so $S = 1$).
% CODE vs DOC: docs/04_launcher.md section 2 writes omega_ball = v/r_c and S = R_b/r_c for the ideal hood;
% src/launcher.py hood_ideal() uses omega = v/R_b (S = 1, rigid). The compressed rolling state S = R_b/r_c
% is what simulate_shot() reaches; both are shown here.

\paragraph{Worked example.} At $n = 3000$\,\unit{rev/min}\assumed{} with no wheel slow-down
($\Omega_f = 314.16$\,\unit{rad/s}\derived), $v = 0.049\times314.16/2 = 7.70$\,\unit{m/s}\derived{} and
$s = v/R_b = 216.8$\,\unit{rad/s}\derived{} (2070\,\unit{rev/min}\derived). With $\delta = 4$\,\unit{mm}\userval,
$S = 0.0355/0.0335 = 1.060$\derived.

\subsection{Equations of motion during contact}
\label{sec:launcher:eom}

\paragraph{Motivation.} Real contacts slip, the wheel slows, and the hood pushes back. The code integrates
these effects through the whole contact rather than assuming the ideal state.

\paragraph{Momentum and angular momentum.} Let $F_w$ and $F_t$ (\unit{N}) be the friction forces \emph{on
the ball} at the wheel and the top contact, positive forward. $M_{w}$ and $M_{t}$ (\unit{N\,m}) are the
rolling-hysteresis moments on the ball (backspin-positive), $\tau_m$ the motor torque at the wheel shaft and
$\tau_b$ the roller bearing drag (both \unit{N\,m}). The code integrates
\begin{align}
  m\,\dot v &= F_w + F_t, \label{eq:launcher-eom_v}\\
  I_b\,\dot s &= r_c\,(F_w - F_t) + M_w + M_t, \label{eq:launcher-eom_s}\\
  I_w\,\dot\Omega &= -R_w F_w + \tau_m(\Omega) + M_w, \label{eq:launcher-eom_W}\\
  I_t\,\dot\Omega_t &= -R_t F_t - \tau_b(\Omega_t) - M_t \quad\text{(roller only)}. \label{eq:launcher-eom_Wt}
\end{align}
Each line is Newton's law for one body. In \eqref{eq:launcher-eom_s} a forward force at the wheel contact
(below the centre) adds backspin and a forward force at the top contact removes it, each with lever arm
$r_c$. In \eqref{eq:launcher-eom_W} the reaction $-F_w$ acts on the wheel rim at radius $R_w$. In the code the
hysteresis reaction on the wheel is \code{-Mw\_wheel} with \code{Mw\_wheel = -Mw\_ball}, which is the
$+M_w$ written here; in \eqref{eq:launcher-eom_Wt}, \code{+Mt\_roll} $= -M_t$.

\paragraph{Normal loads on a concentric hood.} The ball is pressed by the same squeeze at both contacts,
$N_0 = k_n\delta$. Its centre moves on a circle of radius
\begin{equation}
  \rho = R_w + r_c
  \label{eq:launcher-rho}
\end{equation}
about the wheel axis. The wheel pushes the ball outward with $N_w$ and the hood pushes it inward with $N_t$,
so the net inward force must supply the centripetal acceleration and the sum stays at $2k_n\delta$:
\begin{equation}
  N_t - N_w = \frac{m v^2}{\rho}, \qquad N_t + N_w = 2k_n\delta
  \;\;\Longrightarrow\;\;
  N_w = k_n\delta - \frac{m v^2}{2\rho}, \quad N_t = k_n\delta + \frac{m v^2}{2\rho}.
  \label{eq:launcher-Nhood}
\end{equation}
If $N_w$ would become negative the code sets $N_w = 0$, puts the whole centripetal load on the hood,
$N_t = m v^2/\rho$, and raises the flag \code{lost\_wheel\_contact}. The wheel loses contact when
$v^2 > 2\rho k_n\delta/m$. The normal loads use $v$ at the start of each step.

\paragraph{Normal load in a roller nip.} The ball centre travels straight through the gap between the
wheel and the roller. At offset $x$ from the narrowest point, the distance from the wheel axis to the ball
centre is $\sqrt{(R_w+R_b)^2 + x^2} \approx (R_w+R_b) + x^2/[2(R_w+R_b)]$, so the overlap with the wheel
shrinks by $x^2/[2(R_w+R_b)]$; the same holds for the roller with $R_t$. The remaining squeeze and load are
\begin{equation}
  \delta(x) = \delta - \frac{\kappa}{2}x^2, \qquad
  \kappa = \frac{1}{R_w + R_b} + \frac{1}{R_t + R_b}, \qquad
  N_w = N_t = k_n \max\!\big(\delta(x),0\big),
  \label{eq:launcher-Nnip}
\end{equation}
with $\kappa$ in \unit{1/m}. The load vanishes at $\delta(x) = 0$, so the contact runs over
\begin{equation}
  -x_h \le x \le x_h, \qquad x_h = \sqrt{2\delta/\kappa}.
  \label{eq:launcher-xh}
\end{equation}
The roller path has no centripetal term. The hood path runs from $x=0$ to $x_{end} = \theta\rho$
(with $\theta$ in \unit{rad}).

\paragraph{Hysteresis, motor and bearing.} The ball is not a perfect spring. The code applies a rolling
resistance moment $e = c_{rr}\, r_c\, N$ at each contact that opposes the relative rotation of the ball and
the surface, smoothed with $\varepsilon = 1$\,\unit{rad/s}\assumed:
\begin{equation}
  M_w = -c_{rr} r_c N_w \tanh\!\frac{s + \Omega}{\varepsilon}, \qquad
  M_t = -c_{rr} r_c N_t \tanh\!\frac{s - \Omega_t}{\varepsilon}\quad(\Omega_t \equiv 0 \text{ for the hood}).
  \label{eq:launcher-hyst}
\end{equation}
The ball and the wheel turn in opposite senses when the ball has backspin, so their relative rotation is
$s+\Omega$. The motor is a linear DC torque--speed line referred to the wheel shaft (gear ratio $G$ =
motor revolutions per wheel revolution, $V/V_0$ = applied voltage over 12\,\unit{V}\confirmed,
$\omega_{free} = 2\pi n_{free}/60$):
\begin{equation}
  \tau_m(\Omega) = G\,\tau_s\left(\frac{V}{V_0} - \frac{G\,\Omega}{\omega_{free}}\right),
  \label{eq:launcher-motor}
\end{equation}
and the default is $\tau_m = 0$ (the wheel \emph{coasts} through the contact, as in Brettle's analysis).
The roller bearing drag is
\begin{equation}
  \tau_b(\Omega_t) = \tau_c \tanh\!\frac{\Omega_t}{0.5\,\text{rad/s}} + c_v\,\Omega_t .
  \label{eq:launcher-bearing}
\end{equation}

\paragraph{Stick or slip: one step of the solver.} Coulomb's law says $|F_w| \le \mu_w N_w$ and
$|F_t| \le \mu_t N_t$, with equality while the contact slides. The code advances by $\Delta t = 10^{-5}$\,\unit{s}\assumed{}
with forces held constant over the step and chooses $F_w, F_t$ so that the slips \emph{at the end} of the
step obey Coulomb's law. Differentiate \eqref{eq:launcher-slipw} and substitute
\eqref{eq:launcher-eom_v}--\eqref{eq:launcher-eom_W}:
\begin{equation}
  \dot\sigma_w = R_w\dot\Omega - \dot v - r_c\dot s
  = b_w - \Big(\tfrac{R_w^2}{I_w} + \tfrac{1}{m} + \tfrac{r_c^2}{I_b}\Big)F_w
         - \Big(\tfrac{1}{m} - \tfrac{r_c^2}{I_b}\Big)F_t ,
  \label{eq:launcher-sigdot_w}
\end{equation}
where $b_w = R_w\dot\Omega_0 - r_c\dot s_0$ collects the force-free parts
($\dot\Omega_0 = (\tau_m + M_w)/I_w$, $\dot s_0 = (M_w + M_t)/I_b$). The same steps for
\eqref{eq:launcher-slipt} give $\dot\sigma_t = b_t - (\tfrac1m - \tfrac{r_c^2}{I_b})F_w -
(\tfrac{R_t^2}{I_t} + \tfrac1m + \tfrac{r_c^2}{I_b})F_t$ with $b_t = R_t\dot\Omega_{t,0} + r_c\dot s_0$
(the $R_t^2/I_t$ and $R_t\dot\Omega_{t,0}$ terms are zero for a hood). After one step,
\begin{gather}
  \bm{\sigma}^{+} = \Delta t\,\big(\mathbf{q} - \mathbf{A}\mathbf{F}\big),\qquad
  \mathbf{q} = \frac{\bm{\sigma}}{\Delta t} + \mathbf{b}, \nonumber\\
  \mathbf{A} = \begin{pmatrix}
     \frac{R_w^2}{I_w} + \frac1m + \frac{r_c^2}{I_b} & \frac1m - \frac{r_c^2}{I_b}\\[2pt]
     \frac1m - \frac{r_c^2}{I_b} & \frac{R_t^2}{I_t} + \frac1m + \frac{r_c^2}{I_b}
  \end{pmatrix}.
  \label{eq:launcher-lcp}
\end{gather}
Each contact $i$ must be in one of three modes:
\begin{equation}
  \begin{aligned}
  &\text{stick:} && \sigma_i^{+} = 0, && |F_i| \le \mu_i N_i;\\
  &\text{slide forward:} && F_i = +\mu_i N_i, && \sigma_i^{+} \ge 0;\\
  &\text{slide backward:} && F_i = -\mu_i N_i, && \sigma_i^{+} \le 0 .
  \end{aligned}
  \label{eq:launcher-modes}
\end{equation}
$\mathbf{A}$ is symmetric positive definite, so the friction forces that satisfy these conditions are unique;
\code{\_solve\_friction} enumerates the $3\times3 = 9$ mode pairs and keeps the first consistent one. The state is then updated with the constant forces,
$v^{+} = v + \Delta t(F_w + F_t)/m$ and likewise for $s,\Omega,\Omega_t$, and the path position with the
mid-step speed, $x^{+} = x + \Delta t\,(v + v^{+})/2$. The shot ends when $x \ge x_{end}$
(\code{exited}), or is declared stalled if $v \le 0$ after 0.01\,\unit{s}\assumed, or stops at
$t_{max} = 0.25$\,\unit{s}\assumed. Gravity along the path is not modelled.

\paragraph{Energy ledger.} With constant forces in a step, the work done is exactly force times mid-step
velocity, so the code's ledger closes to round-off:
\begin{equation}
  \underbrace{\tfrac12 I_w(\Omega_0^2 - \Omega^2)}_{\text{wheel KE lost}} + W_m
  = \Delta KE_{ball} + \Delta KE_{roller} + E_{f,w} + E_{f,t} + E_{hyst} + E_{bear},
  \label{eq:launcher-ledger}
\end{equation}
where $W_m = \sum \tau_m\bar\Omega\Delta t$, $E_{f,i} = \sum F_i\,\bar\sigma_i\,\Delta t \ge 0$ is the sliding
loss, $E_{hyst} = -\sum[M_w(\bar s + \bar\Omega) + M_t(\bar s - \bar\Omega_t)]\Delta t$ and
$E_{bear} = \sum\tau_b\bar\Omega_t\Delta t$ (bars = mid-step values; all in \unit{J}).

\paragraph{Worked example (nominal hood, 3000\,\unit{rev/min}).} \code{hood\_model(3000)} with the
defaults of Table~\ref{tab:launcher-params}:
$\rho = 0.049 + 0.0335 = 0.0825$\,\unit{m}\derived, path length $\theta\rho = (\pi/2)(0.0825) = 0.1296$\,\unit{m}\derived,
$N_0 = 4000\times0.004 = 16$\,\unit{N}\derived. The ball leaves after 24.6\,\unit{ms}\predicted{} with
$v = 6.996$\,\unit{m/s}\predicted, $s = 208.8$\,\unit{rad/s}\predicted{} (1994\,\unit{rev/min}\predicted),
$S = 1.0597$\predicted{} and both slips zero at exit. The wheel ends at
2726.9\,\unit{rev/min}\predicted{}, and $R_w\Omega_f/2 = 0.049\times285.56/2 = 6.996$\,\unit{m/s}\derived{}
matches $v$, as \eqref{eq:launcher-ideal} requires for a rolling exit.
% CODE vs DOC: docs/04_launcher.md section 3.1 gives the 90 deg hood contact length as ~131 mm; the code
% path is theta*(R_w + r_c) = 129.6 mm (r_c = 33.5 mm). Nip: 2*x_h = 33.4 mm, doc ~34 mm (agrees).

\subsection{Why a fixed hood gives backspin, and the band of $S$}
\label{sec:launcher:Sband}

\paragraph{Motivation.} The hood does not \emph{create} backspin. The wheel does. The hood mainly decides
how much of it survives. Two limiting cases bracket every hood shot.

\paragraph{Limit 1: frictionless top.} Let only the wheel push ($F_t = 0$, $c_{rr} = 0$), with the ball fed
at $v_{in}$ without spin. Dividing \eqref{eq:launcher-eom_s} by \eqref{eq:launcher-eom_v},
\begin{equation}
  \frac{I_b\,\mathrm{d}s}{m\,\mathrm{d}v} = r_c
  \;\;\Longrightarrow\;\;
  s = \frac{m r_c}{I_b}\,(v - v_{in})
  \;\;\Longrightarrow\;\;
  S = \underbrace{\frac{m\,r_c R_b}{I_b}}_{S_\infty}\left(1 - \frac{v_{in}}{v}\right),
  \label{eq:launcher-Sinf}
\end{equation}
and using \eqref{eq:launcher-Ib}, $S_\infty = (r_c/R_b)/k_I$. This is \code{frictionless\_top\_limit}; the code
function returns $S_\infty$, the value for a ball that starts at rest.

\paragraph{Which way does the hood rub?} In the same frictionless case the hood-side point moves at
\begin{equation}
  v - s\,r_c = v - \frac{m r_c^2}{I_b}(v - v_{in}) \approx v\,(1 - \beta), \qquad
  \beta \equiv \frac{m r_c^2}{I_b} = \frac{(r_c/R_b)^2}{k_I},
  \label{eq:launcher-beta}
\end{equation}
for $v \gg v_{in}$. For a hollow ball $\beta > 1$, so the hood-side point moves \emph{backward}, $\sigma_t =
-(v - s r_c) > 0$, and a real hood's friction acts \emph{forward} on the ball ($F_t > 0$). By
\eqref{eq:launcher-eom_s} a forward $F_t$ \emph{removes} backspin and by \eqref{eq:launcher-eom_v} it adds
speed. The hood therefore pulls $S$ down from $S_\infty$ toward the rolling value.
% CODE vs DOC: docs/04 section 3.3 says the hood-side point moves backward "because k_I < 1". In the code
% the condition is beta = m r_c^2 / I_b > 1, i.e. k_I < (r_c/R_b)^2 (= 0.890 at delta = 4 mm, lambda = 1).

\paragraph{Limit 2: rolling on the hood.} If hood friction is high enough or the wrap long enough, the hood
contact sticks and \eqref{eq:launcher-rollhood} holds, so $S = R_b/r_c$ exactly as in
\eqref{eq:launcher-ideal}. For a slow feed, therefore,
\begin{equation}
  \frac{R_b}{r_c} \;\le\; S \;\le\; S_\infty\left(1 - \frac{v_{in}}{v}\right) \;<\; S_\infty .
  \label{eq:launcher-Sband}
\end{equation}
The unit test \code{test\_hood\_S\_bounds} checks $1 \le S\,r_c/R_b \le S_\infty r_c/R_b$ over hood friction
and wrap, for a ball fed from rest ($v_{in} = 0$) with $c_{rr} = 0$.
% CODE vs DOC: docs/04 section 2 item 1 gives the hood-slip band as "between 1 and m r_c R_b / I_b"; the code's
% lower bound is the compressed rolling value R_b/r_c (>1), and the upper bound is reduced by (1 - v_in/v).

\paragraph{Worked example.} With the nominal values, $\beta = 0.8905/0.63 = 1.4135$\derived{} and
$S_\infty = 0.9437/0.63 = 1.498$\derived. Running the code with $\mu_t = 0$, $c_{rr} = 0$ at
3000\,\unit{rev/min}\assumed:
from rest ($v_{in} = 0$) it returns $S = 1.4979$\predicted, equal to $S_\infty$; with
$v_{in} = 0.5$\,\unit{m/s}\assumed{} it returns $v = 6.142$\,\unit{m/s}\predicted{} and
$S = 1.3759$\predicted, and \eqref{eq:launcher-Sinf} gives $1.498\,(1 - 0.5/6.142) = 1.3759$\derived.
The rolling value depends only on the squeeze:

\begin{center}
\small
\begin{tabular}{@{}lccc@{}}
\toprule
$\delta$ & 2\,\unit{mm}\assumed & 4\,\unit{mm}\userval & 8\,\unit{mm}\assumed \\
\midrule
$S = R_b/r_c$ ($\lambda = 1$\assumed) & 1.029\derived & 1.060\derived & 1.127\derived \\
\bottomrule
\end{tabular}
\end{center}

\paragraph{The band in the parameter box.} \code{reachable\_S\_summary} sweeps the fixed hood over
$k_I \in [0.55, 0.667]$\assumed, $\mu_t \in [0.05, 1.0]$\assumed, $\theta \in [45^\circ, 150^\circ]$\assumed,
$\delta \in [2, 8]$\,\unit{mm}\assumed{} and $k_n \in [1.5, 10]$\,\unit{N/mm}\assumed{} at
3000\,\unit{rev/min}\assumed{} (405\derived{} cases). It finds $S$ from 1.029\predicted{} to
1.274\predicted{} and $v$ from 3.18\predicted{} to 7.05\,\unit{m/s}\predicted. A grippy or long hood sits at
the bottom of this band; only a slick, short hood moves toward the top.

\paragraph{$S < 1$ from a fixed hood.} Equation \eqref{eq:launcher-Sband} assumed a slow feed. If the ball
\emph{enters} already moving with little spin and leaves before hood friction rolls it up, the factor
$(1 - v_{in}/v)$ and the unfinished approach to rolling can push $S$ below 1. The stored envelope at
3000\,\unit{rev/min}\assumed{} gives $S = 0.643$\predicted{} for $\theta = 45^\circ$\assumed,
$\mu_t = 0.05$\assumed, $v_{in} = 5$\,\unit{m/s}\assumed. Running the code for $\theta = 20^\circ$\assumed,
$\mu_t = 0.05$\assumed, $v_{in} = 4$\,\unit{m/s}\assumed{} gives $S = 0.575$\predicted.
This needs a fast pre-accelerating feeder and is sensitive to every friction value, so it is not a
design recommendation.
% CODE vs DOC: docs/04 section 3.3 says "20 deg wrap, mu_h 0.05, v_in 4 m/s gives S 0.49". Running
% hood_model(3000, wrap_deg=20, Contact(mu_top=0.05), v_in=4.0) gives S = 0.575 (v = 5.99 m/s);
% S = 0.489 in data/launch_envelope.json belongs to wrap 20 deg, mu_hood 0.10, v_in 5.0 m/s.

\subsection{Rolling and grip-limited regimes}
\label{sec:launcher:grip}

\paragraph{Motivation.} Above some wheel speed, turning the wheel faster stops making the ball faster.
This is the most important practical limit of a single-wheel hood, and it follows directly from Coulomb
friction in \eqref{eq:launcher-modes}.

\paragraph{Wheel sliding, hood sticking.} Suppose the wheel contact slides forward,
$F_w = \mu_w N_w$, while the hood sticks ($v = s r_c$, so $\dot v = r_c \dot s$), with $c_{rr} = 0$. Write
\eqref{eq:launcher-eom_s} as $(I_b/r_c^2)\,\dot v = F_w - F_t$ and add it to \eqref{eq:launcher-eom_v}:
\begin{equation}
  \Big(m + \frac{I_b}{r_c^2}\Big)\dot v = 2F_w
  \;\;\Longrightarrow\;\;
  m_{eff}\,\dot v = F_w, \qquad m_{eff} = \frac{1}{2}\Big(m + \frac{I_b}{r_c^2}\Big).
  \label{eq:launcher-meff}
\end{equation}
Putting $\dot v = F_w/m_{eff}$ back into \eqref{eq:launcher-eom_v} gives the hood force
\begin{equation}
  F_t = \frac{m}{m_{eff}}F_w - F_w = \frac{\beta - 1}{\beta + 1}\,F_w ,
  \label{eq:launcher-Ft}
\end{equation}
which is forward for $\beta > 1$, as found above. The hood keeps sticking only while
$|F_t| \le \mu_t N_t$.

\paragraph{The speed cap.} Along the path $\dot v = v\,\mathrm{d}v/\mathrm{d}x$, so with $w = v^2$ and the
hood load \eqref{eq:launcher-Nhood},
\begin{equation}
  \frac{\mathrm{d}w}{\mathrm{d}x} = \frac{2\mu_w}{m_{eff}}\Big(k_n\delta - \frac{m\,w}{2\rho}\Big)
  = \mathcal{A} - \mathcal{B}w, \qquad
  \mathcal{A} = \frac{2\mu_w k_n\delta}{m_{eff}},\quad \mathcal{B} = \frac{\mu_w m}{m_{eff}\,\rho}.
  \label{eq:launcher-wode}
\end{equation}
This linear equation has the solution, from $w(0) = v_{in}^2$ to $x = \theta\rho$,
\begin{equation}
  v_{cap}^2 = w_\infty - \big(w_\infty - v_{in}^2\big)\,e^{-\mathcal{B}\theta\rho}, \qquad
  w_\infty = \frac{\mathcal{A}}{\mathcal{B}} = \frac{2\rho k_n\delta}{m}, \qquad
  \mathcal{B}\theta\rho = \mu_w\theta\,\frac{2\beta}{1+\beta}.
  \label{eq:launcher-vcap}
\end{equation}
Nothing in \eqref{eq:launcher-vcap} depends on the wheel speed: once the wheel still slips at exit, the exit
speed is fixed by $\mu_w$, $k_n\delta$, the wrap and the ball inertia. $w_\infty$ is exactly the speed at
which the wheel would lose contact. The rpm independence is what \code{test\_grip\_limited\_speed\_independent\_of\_rpm}
checks.

\paragraph{Where the regimes meet.} The wheel catches up with the ball (rolling exit) if
$R_w\Omega_f/2 \le v_{cap}$. Using \eqref{eq:launcher-brettle} from Section~\ref{sec:launcher:drop} for
$\Omega_f$, the crossover wheel speed is
\begin{equation}
  \Omega_i^{*} = a\,\frac{2v_{cap}}{R_w} - \frac{R_w m v_{in}}{2I_w}.
  \label{eq:launcher-crossover}
\end{equation}

\paragraph{Worked example.} Nominal hood: $I_b/r_c^2 = 0.017687$\,\unit{kg}\derived, so
$m_{eff} = 0.021343$\,\unit{kg}\derived; $w_\infty = 2(0.0825)(16)/0.025 = 105.6$\,\unit{m$^2$/s$^2$}\derived{}
($\sqrt{w_\infty} = 10.28$\,\unit{m/s}\derived); $\mathcal{B}\theta\rho = 0.8(\pi/2)(1.1713) = 1.472$\derived;
$e^{-1.472} = 0.2295$\derived. Then $v_{cap}^2 = 105.6 - (105.6 - 0.25)(0.2295) = 81.42$\,\unit{m$^2$/s$^2$}\derived, so
$v_{cap} = 9.02$\,\unit{m/s}\derived. The code with $c_{rr} = 0$ at 5000 and 6000\,\unit{rev/min}\assumed{}
returns 9.021\,\unit{m/s}\predicted{} for both; the small difference is the short initial hood slip that
\eqref{eq:launcher-vcap} ignores. With the default $c_{rr} = 0.02$\assumed{} the stored envelope gives
8.909\,\unit{m/s}\predicted{} from 4000\assumed{} to 6000\,\unit{rev/min}\assumed. The stick check holds:
$(\beta-1)/(\beta+1)\,\mu_w N_w = 0.171\times0.8\,N_w = 0.137\,N_w \le 0.3\,N_t$\derived{} since $N_t \ge N_w$.
Equation \eqref{eq:launcher-crossover} gives $\Omega_i^{*} = 402.7$\,\unit{rad/s}\derived{}
(3845\,\unit{rev/min}\derived{}); the envelope (default $c_{rr}$) is still rolling at 3500\assumed{} and slipping at 4000\,\unit{rev/min}\assumed{}
(predicted regimes).
For a softer ball, $k_n = 1500$\,\unit{N/m}\assumed, the code caps at 5.45\,\unit{m/s}\predicted.

\paragraph{Roller nip cap.} In a nip the load is the parabola \eqref{eq:launcher-Nnip} with no centripetal
term. For a top surface that does not move (a very heavy roller, the same $m_{eff}$), integrating
$\mathrm{d}w/\mathrm{d}x = 2\mu_w N_w/m_{eff}$ over $[-x_h, x_h]$ with
$\int_{-x_h}^{x_h}(\delta - \tfrac{\kappa}{2}x^2)\,\mathrm{d}x = \tfrac43\delta x_h$ gives
\begin{equation}
  v_{cap,nip}^2 = v_{in}^2 + \frac{8\,\mu_w k_n\,\delta\,x_h}{3\,m_{eff}} .
  \label{eq:launcher-vnip}
\end{equation}
With $R_t = 0.024$\,\unit{m}\assumed, $\kappa = 28.64$\,\unit{1/m}\derived{} and
$x_h = 0.01671$\,\unit{m}\derived{} (a 33.4\,\unit{mm}\derived{} contact against 129.6\,\unit{mm}\derived{} for the
hood), \eqref{eq:launcher-vnip} gives $v_{cap,nip} = 5.19$\,\unit{m/s}\derived. The code with a fixed-like
roller and $c_{rr} = 0$ gives 5.18\,\unit{m/s}\predicted; the nominal roller ($I_t = 10^{-5}$\,\unit{kg\,m$^2$}\assumed)
caps at 5.03\,\unit{m/s}\predicted{} from 2500\,\unit{rev/min}\predicted{} upward. The short nip is why a roller is
grip-limited much earlier than a $90^\circ$\assumed{} hood.

\subsection{The ratio $v/u$, and why $v = u/2$ is not universal}
\label{sec:launcher:vu}

\paragraph{Motivation.} Teams set the wheel speed, not the exit speed. The link between them is the
speed ratio
\begin{equation}
  \eta_v \equiv \frac{v}{u_0} = \frac{v}{R_w\Omega_i},
  \label{eq:launcher-etav}
\end{equation}
where $u_0$ and $\Omega_i$ are the surface speed and wheel speed \emph{before} the shot.

\paragraph{Rolling regime.} From \eqref{eq:launcher-ideal} and \eqref{eq:launcher-brettle} (with $c_{rr} = 0$),
\begin{equation}
  \eta_v = \frac{\Omega_f}{2\Omega_i}
         = \frac{1}{2a}\left(1 + \frac{R_w m v_{in}}{2 I_w \Omega_i}\right) < \frac12 ,
  \label{eq:launcher-etav_roll}
\end{equation}
because $a > 1$ and the feed term is small ($R_w m v_{in}/(2I_w) < (a-1)\,\Omega_i$, i.e.\ above about
12\,\unit{rad/s}\derived{} here): the wheel has slowed by the time the ball rolls off it.

\paragraph{Grip-limited regime.} $v = v_{cap}$ is fixed, so $\eta_v = 60\,v_{cap}/(\pi D n)$ falls as
$1/n$.

\paragraph{Other losses.} Hysteresis \eqref{eq:launcher-hyst} lowers $\eta_v$ further. Effects that the code
does \emph{not} include would also lower it: tread compression (an effective wheel radius below $R_w$),
non-diametral contacts, a hood that converges, and frame kick.
% CODE vs DOC: docs/04 section 2 item 2 says the tread also compresses so the effective R_w is R_w - delta_w;
% src/launcher.py keeps R_w fixed and puts all squeeze into the ball radius through lambda.

\paragraph{Worked example.} With $c_{rr} = 0$ at 3000\,\unit{rev/min}\assumed,
\eqref{eq:launcher-etav_roll} gives $(1 + 1.156/314.16)/(2\times1.0967) = 0.4576$\derived; the code gives
0.4576\predicted. With the defaults the code gives the ratios below; at 6000\,\unit{rev/min}\assumed{} the ratio is
barely more than half of the ideal 0.5\derived.

\begin{center}
\small
\begin{tabular}{@{}lccccccc@{}}
\toprule
$n$ [\unit{rev/min}] & 1000\assumed & 2000\assumed & 3000\assumed & 3500\assumed & 4000\assumed & 5000\assumed & 6000\assumed \\
\midrule
$v$ [\unit{m/s}] & 2.239\predicted & 4.640\predicted & 6.996\predicted & 8.169\predicted & 8.909\predicted & 8.909\predicted & 8.909\predicted \\
$\eta_v = v/u_0$ & 0.436\predicted & 0.452\predicted & 0.454\predicted & 0.455\predicted & 0.434\predicted & 0.347\predicted & 0.289\predicted \\
regime & \multicolumn{4}{c}{rolling exit} & \multicolumn{3}{c}{wheel slipping at exit} \\
\bottomrule
\end{tabular}
\end{center}

\paragraph{Literature band.} Because nothing has been measured on POLLEN, the design work uses a band from
\code{research/params\_lit.json} (\code{launcher.v\_over\_surface\_speed}, grade C): low
0.30\assumed, nominal 0.42\assumed, high 0.50\assumed. These come from slow-motion and forum estimates on
other balls and from the rigid ideal; they are literature bounds, not POLLEN measurements. The code adds
0.25\assumed{} (\code{V\_OVER\_U} in \code{src/shooter\_spec.py}) from a multi-roller table-tennis machine
whose measured ratio, about 0.48\assumed{} of the \emph{mean} roller speed, maps to about 0.24--0.25\derived{} of $u_0$ for a
hood (\code{research/EXTRA\_SOURCES.md}). The PM-001 model's own nominal hood, 0.436--0.455\predicted{} below the grip
limit, sits on the optimistic side of that band.
% NOTE: src/shooter_spec.py comments the 0.25 value as "Sakai & Kitade 2017 low end"; research/EXTRA_SOURCES.md
% (row launcher.v_over_surface_speed) says it is a CALC mapping of eta ~0.48-0.51 (relative to mean roller
% surface speed) to v/u_w for a hood, not a measured low end.

\subsection{Flywheel slow-down and energy per shot}
\label{sec:launcher:drop}

\paragraph{Motivation.} A coasting wheel gives up speed to the ball. That drop sets $\Omega_f$ in
\eqref{eq:launcher-ideal}, the recovery time between shots, and how much of the motor's work reaches the ball.

\paragraph{Impulse bookkeeping.} Integrate \eqref{eq:launcher-eom_v}--\eqref{eq:launcher-eom_W} over the
contact with $\tau_m = 0$, $c_{rr} = 0$ and define the friction impulses
$J_w = \int F_w\,\mathrm{d}t$, $J_t = \int F_t\,\mathrm{d}t$ (\unit{N\,s}):
\begin{equation}
  m(v - v_{in}) = J_w + J_t, \qquad \frac{I_b s}{r_c} = J_w - J_t, \qquad
  I_w(\Omega_i - \Omega_f) = R_w J_w .
  \label{eq:launcher-impulse}
\end{equation}
Adding the first two eliminates the unknown hood impulse:
\begin{equation}
  m(v - v_{in}) + \frac{I_b s}{r_c} + \frac{2 I_w(\Omega_f - \Omega_i)}{R_w} = 0 ,
  \label{eq:launcher-momhood}
\end{equation}
which \code{test\_generalised\_momentum\_hood} checks to $10^{-9}$.

\paragraph{Generalised Brettle drop.} For a rolling exit, $v = s r_c = R_w\Omega_f/2$. Then
$J_w = \tfrac12[m(v - v_{in}) + I_b v/r_c^2]$, and the wheel equation in \eqref{eq:launcher-impulse} becomes
\begin{equation}
  I_w\Omega_i - I_w\Omega_f = \frac{R_w}{2}\left[\Big(m + \frac{I_b}{r_c^2}\Big)\frac{R_w\Omega_f}{2} - m v_{in}\right].
  \label{eq:launcher-brettlestep}
\end{equation}
Collecting the $\Omega_f$ terms and dividing by $I_w$:
\begin{equation}
  \Omega_f = \frac{\Omega_i + \dfrac{R_w m v_{in}}{2 I_w}}{a}, \qquad
  a = 1 + \frac{R_w^2\,(m + I_b/r_c^2)}{4 I_w},
  \label{eq:launcher-brettle}
\end{equation}
which is \code{brettle\_omega\_f} and \code{brettle\_ratio}. For $r_c = R_b$ and $v_{in} = 0$, using
\eqref{eq:launcher-Ib}, $a = 1 + (1 + k_I)\,m R_w^2/(4 I_w)$, Brettle's published ratio. The drop per shot is
$\Delta n = n_i - n_f$.

\paragraph{Grip-limited drop.} If the wheel still slips at exit, the hood still sticks, so $s = v_{cap}/r_c$
and the same $J_w$ applies with $v = v_{cap}$:
\begin{equation}
  \Delta\Omega = \frac{R_w}{2 I_w}\Big[m(v_{cap} - v_{in}) + \frac{I_b\,v_{cap}}{r_c^2}\Big],
  \label{eq:launcher-dropgrip}
\end{equation}
which is also independent of wheel speed.

\paragraph{Efficiency.} With $v_{in} = 0$, the ball's kinetic energy at a rolling exit is
$\tfrac12(m + I_b/r_c^2)v^2 = \tfrac12 I_w\Omega_f^2\,(a - 1)$ (substitute $v = R_w\Omega_f/2$ and the
definition of $a$). The wheel loses $\tfrac12 I_w(\Omega_i^2 - \Omega_f^2) = \tfrac12 I_w\Omega_f^2(a^2 - 1)$.
Their ratio is
\begin{equation}
  \eta = \frac{KE_{ball}}{\Delta KE_{wheel}} = \frac{a - 1}{a^2 - 1} = \frac{1}{a + 1} < \frac12 .
  \label{eq:launcher-eta}
\end{equation}
The rest, $a$ times the ball's energy, is lost to sliding at the contacts while the ball spins up (in the
nominal run below almost all of it at the wheel).

\paragraph{Recovery.} After the shot the motor \eqref{eq:launcher-motor} at full voltage refills the wheel:
$I_w\dot\Omega = T_0(1 - \Omega/\Omega_{free})$ with $T_0 = G\tau_s V/V_0$ and
$\Omega_{free} = (V/V_0)\,\omega_{free}/G$. Separating variables,
\begin{equation}
  t_{rec} = \tau_{mw}\ln\frac{\Omega_{free} - \Omega_{start}}{\Omega_{free} - \Omega_{target}}, \qquad
  \tau_{mw} = \frac{I_w\Omega_{free}}{T_0},
  \label{eq:launcher-recovery}
\end{equation}
infinite when $\Omega_{target} \ge \Omega_{free}$ (\code{recovery\_time}; no current limit, no controller lag,
no rotor inertia).

\paragraph{Worked example.} Rigid Brettle case: $a = 1 + 0.049^2(0.025 + 0.01575)/(4\times2.65\times10^{-4}) =
1.0923$\derived, a drop of $1 - 1/a = 8.45$\,\%\derived{}, i.e.\ 253.5\,\unit{rev/min}\derived{} at
3000\,\unit{rev/min}\assumed{} (the \code{hood\_ideal\_brettle} row). With $r_c = 0.0335$\,\unit{m}\derived,
$a = 1.0967$\derived; adding $v_{in} = 0.5$\,\unit{m/s}\assumed, \eqref{eq:launcher-brettle} gives
$n_f = 2745.6$\,\unit{rev/min}\derived{}, a drop of 254.4\,\unit{rev/min}\derived, and the code with
$c_{rr} = 0$ gives 254.43\,\unit{rev/min}\predicted. With hysteresis the nominal drop is
273.1\,\unit{rev/min}\predicted. In the grip-limited regime \eqref{eq:launcher-dropgrip} with
$v_{cap} = 9.021$\,\unit{m/s}\predicted{} gives 328.9\,\unit{rev/min}\derived{} (code, $c_{rr}=0$:
328.9\,\unit{rev/min}\predicted; default: 343.6\,\unit{rev/min}\predicted).
Equation \eqref{eq:launcher-eta} gives $\eta = 1/2.0967 = 0.477$\derived. The nominal ledger at
3000\,\unit{rev/min}: wheel loses 2.273\,\unit{J}\predicted, the ball gains 1.042\,\unit{J}\predicted{}
($\eta = 0.458$\predicted), wheel-side sliding 1.093\,\unit{J}\predicted, hood sliding
0.0009\,\unit{J}\predicted, hysteresis 0.137\,\unit{J}\predicted; the residual is zero to the printed digits.
Recovery with the goBILDA 5203 1:1\confirmed{} at 12\,\unit{V}\confirmed:$\Omega_{free} = 628.3$\,\unit{rad/s}\derived,
$\tau_{mw} = 2.65\times10^{-4}\times628.3/0.144 = 1.156$\,\unit{s}\derived,
$\Omega_{start} = 314.16 - 28.60 = 285.56$\,\unit{rad/s}\derived, so
$t_{rec} = 1.156\ln(342.76/314.16) = 0.101$\,\unit{s}\derived. Holding full voltage \emph{during} contact
instead of coasting gives 7.153\,\unit{m/s}\predicted{} instead of 6.996\,\unit{m/s}\predicted{} at
3000\,\unit{rev/min}\assumed, and adds nothing in the grip-limited regime.

\subsection{The passive upper roller}
\label{sec:launcher:roller}

\paragraph{Motivation.} A common idea is that a free top roller ``spins up and reduces the spin
difference''. The model tests that with the roller equation \eqref{eq:launcher-eom_Wt}.

\paragraph{Rolling state with a moving top surface.} Let $u_w = R_w\Omega$ and $u_t = R_t\Omega_t$ be the
two surface speeds (forward positive). Rolling at both contacts means $\sigma_w = \sigma_t = 0$:
$u_w = v + s r_c$ and $u_t = v - s r_c$. Adding and subtracting,
\begin{equation}
  v = \frac{u_w + u_t}{2}, \qquad s = \frac{u_w - u_t}{2 r_c}, \qquad
  S = \frac{R_b}{r_c}\cdot\frac{u_w - u_t}{u_w + u_t},
  \label{eq:launcher-rollerstate}
\end{equation}
which is \code{roller\_rolling\_state}. A fixed hood is $u_t = 0$. $S$ drops below $R_b/r_c$ only if the top
surface moves \emph{forward} ($u_t > 0$).

\paragraph{Which way the passive roller turns.} By \eqref{eq:launcher-beta}, for a slow feed the ball's top
point moves backward, so it drags the roller to $\Omega_t < 0$, i.e.\ $u_t < 0$, and
\eqref{eq:launcher-rollerstate} gives $S > R_b/r_c$. This needs $v \gg v_{in}$: at
1000\,\unit{rev/min}\assumed{} the 0.5\,\unit{m/s}\assumed{} feed is a large part of the exit speed, the nominal
roller is spun \emph{forward} ($+8.25$\,\unit{rad/s}\predicted) and $S = 0.974$\predicted; the backward drag
appears from about 2000\,\unit{rev/min}\predicted{} upward. The two inertia limits are: a very heavy roller at rest
behaves like a fixed hood ($S \to R_b/r_c$), and a massless frictionless roller transmits no force, giving
\eqref{eq:launcher-Sinf} (checked by \code{test\_massless\_roller\_equals\_frictionless\_top}).

\paragraph{Between shots.} The roller keeps its exit speed and coasts for $\Delta t_{shot}$ under
$I_t\dot\Omega_t = -\tau_b(\Omega_t)$ (\code{roller\_coast}, 1\,\unit{ms}\assumed{} explicit steps, set to zero
if the sign would flip).

\paragraph{Worked example.} At 2000\,\unit{rev/min}\assumed{} the nominal roller exits rolling at both
contacts with $u_w = 9.375$\,\unit{m/s}\predicted{} and $u_t = -0.166$\,\unit{m/s}\predicted. Equation
\eqref{eq:launcher-rollerstate} gives $v = 4.605$\,\unit{m/s}\derived{} and
$S = 1.0597\times9.541/9.210 = 1.098$\derived, matching the simulated $S = 1.0978$\predicted. At
3000\,\unit{rev/min} it is grip-limited: $v = 5.029$\,\unit{m/s}\predicted, $S = 1.107$\predicted, roller at
$-89.9$\,\unit{rev/min}\predicted{} ($-9.41$\,\unit{rad/s}\predicted). Coasting: the bearing torque is about
$5\times10^{-4} + 10^{-6}\times9.41 = 5.09\times10^{-4}$\,\unit{N\,m}\derived, a deceleration of
50.9\,\unit{rad/s$^2$}\derived, so the roller stops in about 0.19\,\unit{s}\derived{} $< 0.5$\,\unit{s}\assumed.
Every shot of a burst therefore starts from rest, and the stored \code{roller\_nominal\_steady} rows equal
\code{roller\_nominal\_rest}.

\paragraph{Roller band.} Over $I_t \in [10^{-6}, 2\times10^{-4}]$\,\unit{kg\,m$^2$}\assumed,
$R_t \in [12.5, 50]$\,\unit{mm}\assumed, bearing torque 0--$3\times10^{-3}$\,\unit{N\,m}\assumed{} and
$\delta \in \{4, 8\}$\,\unit{mm}\assumed{} (shots 1--3 of a burst, 216\derived{} cases) the code finds
$S$ from 1.060\predicted{} to 1.242\predicted{} and $v$ from 4.86\predicted{} to 7.09\,\unit{m/s}\predicted{} at
3000\,\unit{rev/min}\assumed. With a fast feed the roller is spun forward: a light roller
($I_t = 2\times10^{-6}$\,\unit{kg\,m$^2$}\assumed) gives $S = 0.921$\predicted{} at $v_{in} = 2$\,\unit{m/s}\assumed{}
and 0.674\predicted{} at 4\,\unit{m/s}\assumed. So the passive roller, fed slowly at a working wheel speed, slightly \emph{raises} $S$ and
costs exit speed; it does not reduce spin.
% NOTE: research/params_lit.json launcher.S_exit_passive_roller.vs_PM001 quotes the PM-001 envelope as
% "0.91-1.24"; data/launch_envelope.json reachable_S_summary.roller gives 1.06-1.242 (slow feed). The 0.91
% value is the fast-feed case (v_in = 2 m/s), listed separately.

\subsection{Inverse problem: wheel speed for a target exit speed}
\label{sec:launcher:inverse}

\paragraph{Motivation.} The trajectory sections give a target exit speed $v$. The team needs a wheel
speed. Inverting \eqref{eq:launcher-etav} with \eqref{eq:launcher-u}:
\begin{equation}
  u_0 = \frac{v}{\eta_v}
  \;\;\Longrightarrow\;\;
  n = \frac{60\,u_0}{\pi D} = \frac{60\,v}{\eta_v\,\pi D},
  \label{eq:launcher-inverse}
\end{equation}
which is \code{wheel\_rpm} in \code{src/shooter\_spec.py}. It is rounded to a whole rpm and
evaluated at
$\eta_v \in \{0.25, 0.30, 0.42, 0.50\}$\assumed. The matching ball spin for a design $S$ follows from
\eqref{eq:launcher-S}: $n_{ball} = |S|\,v\,60/(2\pi R_b)$ (\code{spin\_rpm}).

\paragraph{Worked example.} For the exit-speed band 6.5--9.0\,\unit{m/s}\predicted{} of the shooter
specification:

\begin{center}
\small
\begin{tabular}{@{}lcccc@{}}
\toprule
target $v$ & $\eta_v = 0.25$\assumed & $0.30$\assumed & $0.42$\assumed & $0.50$\assumed \\
\midrule
6.5\,\unit{m/s}\predicted & 5067\,\unit{rev/min}\derived & 4222\derived & 3016\derived & 2533\derived \\
7.0\,\unit{m/s}\predicted & 5457\derived & 4547\derived & 3248\derived & 2728\derived \\
9.0\,\unit{m/s}\predicted & 7016\derived & 5847\derived & 4176\derived & 3508\derived \\
\bottomrule
\end{tabular}
\end{center}

For example $n = 60\times9.0/(0.25\,\pi\times0.098) = 7016$\,\unit{rev/min}\derived. Two checks against the
forward model. First, 7016\,\unit{rev/min}\derived{} is above the 6000\,\unit{rev/min}\confirmed{} free speed of a 1:1
motor, so the 0.25 case needs a step-up drive or a different motor. Second, the inverse
\eqref{eq:launcher-inverse} assumes a constant $\eta_v$, which only holds in the rolling regime. The nominal
hood of Section~\ref{sec:launcher:grip} cannot reach 9.0\,\unit{m/s}\predicted{} at \emph{any} wheel speed (cap
8.91\,\unit{m/s}\predicted); a higher cap needs more grip, more squeeze or a stiffer ball, which is what
measuring $\mu_w$, $k_n$ and $\eta_v$ is for. Spin for $S = 1$ at 7.0\,\unit{m/s}\predicted: 1883\,\unit{rev/min}\derived.
% CODE vs DOC: docs/08_shooter_spec.md quotes 3020-4180 / 4220-5850 / 5070-7020 / 2530-3510 rpm for
% 6.5-9.0 m/s; shooter_spec.wheel_rpm gives 3016-4176 / 4222-5847 / 5067-7016 / 2533-3508 (doc rounded to 10 rpm).

\subsection{Assumptions}
\label{sec:launcher:assumptions}

\begin{itemize}
  \item Planar, one-dimensional contact path: a hood concentric with the wheel (uniform gap) or a straight
        pass through a two-cylinder nip. No converging hood, entry ramp, off-axis contact, sidespin or frame
        kick.
  \item Linear normal law $N = k_n\delta$ with an unmeasured $k_n$; a real hollow shell is non-linear and can
        buckle. The squeeze shortens the rolling radius by $\lambda\delta/2$ with unknown $\lambda$; the tread is
        rigid and $R_w$ is fixed.
  \item Constant Coulomb coefficients $\mu_w$, $\mu_t$ (no dependence on slip speed, heat or wear); none is
        measured on POLLEN.
  \item Rolling hysteresis as a moment $c_{rr} r_c N$ smoothed with $\tanh$; $c_{rr}$ is a guess.
  \item The ball enters with speed $v_{in}$ and no spin; gravity along the path is ignored.
  \item The wheel coasts during contact unless a motor is given; the motor is a linear torque--speed line
        with no current limit, controller lag or rotor inertia.
  \item Ball mass, radius and inertia vary between pieces only through the parameter boxes; $k_I$ is derived
        from a uniform shell, not measured.
  \item The speed ratios $\eta_v$ used in the inverse problem are literature bounds from other balls and
        launchers, not POLLEN measurements.
  \item Every simulated value in this section is a prediction of an uncalibrated model. A single test shot
        would not establish repeatability; the calibration needs repeated shots at several wheel speeds.
\end{itemize}
